\chapter*{Introduction}%
\addcontentsline{itc}{chapter}{Introduction}

\section*{Overview of the monograph}

This monograph is devoted to the study
of Hilbertian correlations (also called ``maximal correlations'' or ``$\rho$-mixing coefficients''),
in particular to showing how this concept can be `tensorized'
to yield new results on systems of statistical mechanics exhibiting asymptotic independence.
I have divided it into six chapters:
%
\begin{itemize}
\item The first chapter, numbered ``\ref{parMotivation}'',
aims at motivating the study of Hilbertian correlations and their tensorization.
In this chapter, I will recall some classical results on the subcritical Ising model,
which is a classical model showing asymptotic independence between pairs of spins.
When one gets interested in very large `bunches' of spins,
it is known that asymptotic independence cannot be captured by $\beta$-mixing any more,
but that, in certain cases at least, it still holds in terms of $\rho$-mixing.
The techniques used so far to establish $\rho$-mixing for bunches of spins
are strongly limited by technical assumptions looking somehow artificial,
which will motivate studying $\rho$-mixing `for itself' and trying to tensorize it.
%
\item In Chapter~\ref{parBasicHilbertianDecorrelations},
I shall recall the definition of the Hilbertian correlation coefficient;
I shall also recall some classical facts about this concept and give some examples.
This chapter can be seen as a `crash course' on $\rho$-mixing for the non-specialist reader:
almost nothing in it is new.
%
\item In Chapter~\ref{parEventSufficientConditions}, I shall give some new criteria
to bound the Hilbertian correlation between two $\sigma$-algebras,
which criteria assume bounds on the~$(\Pr[A\cap B] - \Pr[A]\Pr[B])$
for events~$A$ and~$B$ belonging to these respective $\sigma$-algebras.
My ``strong event sufficient condition'', which improves previous results by several authors,
shall even be shown to be optimal.
\item Chapter~\ref{parTensorization} is the core of this work:
in it I will handle tensorization of Hilbertian decorrelations.
This chapter begins with a refined version of the concept of correlation,
called ``subjective correlation'', which is necessary to write the subsequent tensorization results.
Then I shall state and prove my three main tensorization theorems:
Theorem~\ref{thm5750} (`$N$~against~$1$' theorem)
bounds the correlation between
a `simple' and a `vector' variable;
Theorem~\ref{thm6413} (`$N$~against~$M$' theorem)
deals with correlation between two vector variables,
and Theorem~\ref{thm0119} (`$\ZZ$~against~$\ZZ$' theorem)
refines the previous one in the case where certain symmetries are present.
%
Then, I will discuss some refinement and optimality statements
about these theorems; in \S~\ref{par3lines},
I will also present a geometric corollary of tensorization results
which underlines quite well the Hilbertian aspect of maximal correlations.
%
\item In Chapter~\ref{parOtherApplicationsOfTensorizationTechniques},
I will continue to use the tensorization techniques of Chapter~\ref{parTensorization},
but that time instead of proving tensorization results \emph{stricto sensu}
I will turn to different types of results, namely the spatial central limit theorem
and the presence of spectral gap for the Glauber dynamics.
%
\item Finally, Chapter~\ref{parApplications} will present some concrete applications
of the results of this monograph. For instance, I shall prove new results about decorrelation
between distant bunches of spins in Ising's model (see Theorem~\ref{t3730});
I will also give results of the same type for quite general models of statistical mechanics (see e.g.\ Theorems~\ref{thmTchou} and~\ref{thmLama}),
also proving spatial CLT and spectral gap for the Glauber dynamics for these models.
I will also show how tensorization of Hilbertian correlations can be used
to get `hypocoercivity' results~[Theorem~\ref{thmParis}].
\end{itemize}


\section*{Conventions and notation}

Notation will not always be perfectly rigorous:
to make reading easier, it may occur sometimes that formalism is slightly loose,
or that some writing conventions or assumptions are implicit.
However this shall only be done in situations
where adding the missing information by the reader is (hopefully) obvious.

Here is some notation used throughout this text:

\paragraph{Miscellaneous}

\begin{itemize}
\item The symbol $\NN$ denotes the set of nonnegative integers, including $0$.
The set of positive integers $\NN\setminus\{0\}$ is denoted by~$\NN^*$.
%
\item For~$a,b$ real numbers, $a\wedge b$ denotes $\min\{a,b\}$,
resp.~$a\vee b$ denotes $\max\{a,b\}$; $a_+$ denotes the positive part of~$a$,
i.e.\ $a\vee0$.
%
\item For~$A$ a set, $\C{A}$ denotes the complement set of~$A$
(the set of reference shall always be clear);
$\1{A}$ denotes the indicator function of~$A$,
that is, the function being~$1$ on~$A$ and~$0$ on~$\C{A}$.
%
\item For~$A,B$ sets, $A\bigtriangleup B$ denotes the symmetric difference of~$A$ and~$B$,
i.e.\ ${(A\setminus B)}\uplus(B\setminus A)$,
where ``$\uplus$'' means the same as ``$\cup$'', but with underlining that the union is disjoint.
%
\item The identity matrix in dimension $n$ will be denoted by~$\mathbf{I}_n$.
The transpose of a matrix $A$ will be denoted by~$\T{A}$.
%
\item If $\Theta$ is a set endowed with a metric $\dist$, then for
$I,J\subset\Theta$, $\dist(I,J)$ denotes the distance between~$I$ and~$J$, that is,
$\dist(I,J) \coloneqq \inf\{\dist(i,j) \colon {i\in I}, \ab {j\in J}\}$.
%
\item As is customary in physical literature, $\propto$ means ``proportional to''.
%
\item Whenever $I$ is a set and $X$ a symbol,
$\vec{X}_I$ will be a shorthand for ``$(X_i)_{i\in I}$''.
\end{itemize}


\paragraph{Probability}

\begin{itemize}
\item We will always work on an implicit probability space $(\Omega,\mcal{B})$
equipped with a probability measure~$\Pr$. Sub-$\sigma$-algebras of~$\mcal{B}$
will be merely called ``$\sigma$-algebras'';
I will also often write ``variable'' for ``random variable''.
Unless explicitly specified, variables on~$\Omega$ can be valued in any set.
%
\item If $f$ is a real random variable, the expectation of~$f$ is denoted by~$\EE[f]$;
its variance is denoted by~$\Var(f)$; its standard deviation is denoted
by $\ecty(f) \coloneqq \sqrt{\Var(f)}$; if $g$ is another real variable,
the covariance between~$f$ and~$g$ is denoted by~$\Cov(f,g) \coloneqq \EE[fg] - \EE[f]\EE[g]$.
All that notation extends to the case where $f$ and~$g$ are valued
in some vector space $\RR^N$, except that in that case it refers to vectors or matrices.
%
\item If $B$ is an event with $\Pr[B]>0$, then $\Pr[A|B]$, $\EE[f|B]$, $\Var(f|B)$, \dots\ 
stand resp.\ for the probability of~$A$, the expectation of~$f$, the variance of~$f$, \dots\ 
under the conditional law~$\dx\Pr[\Bcdot|B] \coloneqq \1{B}\,\dx\Pr[\Bcdot] \div\Pr[B]$.
Similarly, if $\mcal{F}$ is a $\sigma$-algebra,
$\Pr[A|\mcal{F}]$, $\EE[f|\mcal{F}]$, \dots\ stand for the conditional probability of~$A$,
the conditional expectation of~$f$, \dots\ w.r.t.~$\mcal{F}$.
%
\item Concerning conditional expectations, I will actually use two different conventions:
for~$\mcal{G}$ a $\sigma$-algebra, $\EE[f|\mcal{G}]$ can also be denoted by~$f^{\mcal{G}}$.
Both conventions can be used inside the same formula.%
\footnote{The use of the first or the second convention will depend on
the way we prefer to see the conditional expectation of~$f$ w.r.t.~$\mcal{G}$:
if it is rather seen as the expectation of~$f$ knowing the information of~$\mcal{G}$,
notation $\EE[f|\mcal{G}]$ will be chosen, while if it is more seen like
the $\mcal{G}$-measurable function best approximating $f$, we will use the notation $f^{\mcal{G}}$.}%
\footnote{One must not confuse $\Var(f|\mcal{G})$,
which is the variance of~$f$ under the law $\Pr[\Bcdot|\mcal{G}]$,
with~$\Var(f^{\mcal{G}})$ which is the (unconditioned) variance of the random variable $f^{\mcal{G}}$.
One has the well-known identity $\Var(f) = \Var(f^{\mcal{G}}) + \EE[\Var(f|\mcal{G})]$,
which I shall refer to as \emph{associativity of variance}.}
%
\item If $X$ is a variable on~$\Omega$, the $\sigma$-algebra generated by~$X$
(that is, the smallest $\sigma$-algebra w.r.t.\ which $X$ is measurable)
is denoted by~$\sigma(X)$. If $\mcal{F}$ and~$\mcal{G}$ are $\sigma$-algebras,
the $\sigma$-algebra generated by~$\mcal{F}$ and~$\mcal{G}$
(that is, the smallest $\sigma$-algebra containing both~$\mcal{F}$ and~$\mcal{G}$)
is denoted by~$\mcal{F} \vee \mcal{G}$,
and for an arbitrary number of $\sigma$-algebras this notation extends into the $\infty$-ary operator~$\bigvee$.
%
\item An event $A\in\mcal{B}$ is said to have trivial probability, or to be trivial,
if $\Pr[A]\in\{0,1\}$. A $\sigma$-algebra is said to be trivial if all its events are trivial.
The $\sigma$-algebra $\{\emptyset,\Omega\}$, which is trivial under any law $\Pr$,
will be denoted by~$\mcal{O}$ and refered to as ``the'' trivial sigma-algebra.
%
\item The Lebesgue measure on~$\RR^n$ will be denoted by~$\dx{x}$,
``$x$'' being the name of the integration variable.
For a Borel set $A\subset\RR^n$, $\int_{x\in A} \dx{x}$ will sometimes be denoted by~$|A|$.
%
\item For~$C$ a positive-semidefinite matrix (possibly of dimension $1$,
in which case it is identified with~$\sigma^2 \in \RR_+$), $\mcal{N}(C)$ denotes
the law of the centered Gaussian vector with covariance matrix~$C$.
I will write $\mcal{N}(C)+m$
to denote the non-centered Gaussian vector with variance $C$ and mean $m$.
\end{itemize}


\paragraph{Functional analysis}

\begin{itemize}
\item Unless otherwise specified, all the functional spaces considered in this monograph
shall be real.
%
\item For~$I$ an open interval of~$\RR$ and~$k\in\NN\cup\{\infty\}$,
$\mcal{C}^k_0(I)$ denotes the subset of functions of~$\mcal{C}^k(I)$ with compact support.
%
\item If $\mu$ is a nonnegative measure on some measurable space $(\Omega,\mcal{B})$,
$L^2(\mu)$ denotes the set of measurable functions $f$ (up to~$\mu$-a.e.\ equality) such that
$\int_{\Omega} f(\omega)^2 \,\dx\mu(\omega) < \infty$.
If $I$ is a countable set, $L^2(I)$ denotes the set of functions $f : I\to\RR$ such that
$\sum_{i\in I} f(i)^2 < \infty$.
If $\mcal{F}$ is a $\sigma$-algebra, $L^2(\mcal{F})$ denotes the space of
$\mcal{F}$-measurable functions (up to a.s.\ equality)
which are square-integrable w.r.t.~$\Pr$.
All these spaces are equipped with their natural Hilbertian product
$\langle\Bcdot,\Bcdot\rangle$ and the associated norm $\|\Bcdot\|$.
%
\item For~$\mu$ a finite measure, in~$L^2(\mu)$ the constant functions make a line
which can be identified with~$\RR$;
then, $\ldb(\mu)$ will denote the quotient $L^2(\mu)/\RR$,
equipped with its natural Hilbert structure.
In other words, if $\bar{f}\in \ldb(\mu)$ is the projection of~$f\in L^2(\mu)$,
$\|\bar{f}\|_{\ldb} \coloneqq \inf\{\|f-a\|_{L^2}\colon {a\in\RR}\} =
\big( \|f\|_{L^2}^2 - \langle f,1/\|1\|\rangle_{L^2}^2 \big)^{1/2}$.
$\ldb(\mu)$ can also be seen as the subspace of centered functions of~$L^2(\mu)$,
i.e.\ as $\{ f\in L^2(\mu)\colon \langle f,1\rangle = 0 \}$;
throughout the monograph we will implicitly switch between both interpretations.
%
\item If $L \colon H_1\to H_2$ is a linear operator between two Hilbert spaces,
then $L^* \colon H_2\to H_1$ denotes the adjoint operator of~$L$, characterized by the relationship
$\langle L^*y, x \rangle_{H_1} = \langle y, Lx \rangle_{H_2}$.
%
\item If $L \colon E\to F$ is a linear operator between two Banach spaces
(not necessarily Hilbert) with respective norms~$\|\Bcdot\|_E$ and~$\|\Bcdot\|_F$,
the operator norm of~$f$, denoted by~$\VERT f\VERT$, is defined as
$\sup\{\|Lx\|_F\colon \|x\|_E=1\}$.
%
\item If $L \colon E\to E$ is a bounded linear operator on a Banach space,
then $\rho(L)$ denotes the spectral radius of~$f$, that is,
$\rho(L) \coloneqq \lim_{k\longto\infty} \VERT L^k \VERT^{1/k}$
—this limit always exists.
%
\item A column vector $(a_i)_{i\in I}$ will automatically be identified with
the corresponding element of~$L^2(I)$.
Likewise, a matrix $A = (\!(a_{ij})\!)_{(i,j)\in I\times J}$ will be identified
with the corresponding linear operator from $L^2(J)$ to~$L^2(I)$.
%
\item In our computations we will often use the Cauchy\,-\,Schwarz inequality
and its variants%
\footnote{For instance, the discrete form $\big|\sum_{i=1}^N a_ib_i\big|
\leq \big(\sum_{i=1}^N a_i^2\big)^{1/2}\big(\sum_{i=1}^N b_i^2\big)^{1/2}$,
the probabilistic form $|\Cov[fg]| \leq \ecty(f)\ecty(g)$, etc..};
when using such an inequality,
we will indicate it by writing ``CS'' under the inequality sign concerned.
Similarly, ``IP'' under an equality sign will mean
that this equality follows from integrating by parts.
\end{itemize}




\section*{Acknowledgements}

This work owes to many people's help.
First of all I must mention V.~Beffara, who launched me on this topic incidentally.
His relevant comments, as well of those of (among others)
Y.~Ollivier, C.~Villani, Y.~Velenik, T.~Bodineau, C.~Shalizi and R.~Bradley,
have also been the origin for several important improvements of this monograph.

Several colleagues provided me with some help on mathematical topics
I was not familiar with. In particular,
Y.~Ollivier suggested to me the use of Lipschitz spaces to prove Lemma~\ref{lem7792};
S.~Martineau pointed out how Gershgorin's Lemma solved a technical point
in the complete proof of Theorem~\ref{thm6657-0};
V.~Calvez had the idea of using Laplace transform to prove Lemma~\ref{l1939a}.
Over the Internet, F.~Martinelli and S.~Shlosman also gave me
precious bibliographic references on the state of the art
about weak and strong mixing in statistical mechanics.

Most of the drawings in this monograph were made thanks to the excellent \LaTeX\ extension \hbox{Ti\textit{k}Z}, combined with computations in C language. The dice of Figure~\ref{fig-3lines} have been kindly drawn for me by A.~Alvarez, using POV-Ray.

